\documentclass[11pt, a4paper]{article}
\usepackage[T1]{fontenc}
\usepackage[utf8]{inputenc}
\usepackage{lmodern}
\usepackage[margin=22mm]{geometry}
\usepackage{amsmath}
\usepackage{enumitem}
\usepackage{tcolorbox}
\usepackage[table]{xcolor}
\usepackage{booktabs}
\usepackage[hidelinks]{hyperref}

\definecolor{accent}{RGB}{43,108,176}
\definecolor{warn}{RGB}{192,72,40}
\newtcolorbox{keybox}[1]{colback=blue!4, colframe=accent, fonttitle=\bfseries,
  title=#1, boxrule=0.6pt, arc=2pt, left=3mm, right=3mm, top=2mm, bottom=2mm}
\newtcolorbox{warnbox}[1]{colback=red!3, colframe=warn, fonttitle=\bfseries,
  title=#1, boxrule=0.6pt, arc=2pt, left=3mm, right=3mm, top=2mm, bottom=2mm}

\newcommand{\answerspace}[1]{\vspace{#1}\noindent\rule{\linewidth}{0.3pt}\\[2mm]
  \noindent\rule{\linewidth}{0.3pt}\\[2mm]\noindent\rule{\linewidth}{0.3pt}\vspace{3mm}}

\sloppy
\title{\vspace{-14mm}\textbf{Sizing Lab Worksheet}\\[2mm]
  \large F-16A, the two-state sizing loop}
\author{AOE 4065 --- Air Vehicle Design}
\date{}

\begin{document}
\maketitle
\vspace{-6mm}

\begin{keybox}{Before you start}
Finish \texttt{lab/LAB01\_write\_the\_L2\_loop.mlx} and get a \textbf{PASS}
from the checker. Every question below is an experiment you run on the loop
you wrote, so you need a working one first.

Work in a copy of your file. Put the setting back after each experiment,
otherwise the next one starts from the last one.
\end{keybox}

% ------------------------------------------------------------------
\section*{Part A --- Under-relaxation}

Your loop ends each pass with
\[
W \leftarrow W + w_W\,(W^{new} - W), \qquad
T \leftarrow T + w_T\,(T^{new} - T).
\]

\begin{enumerate}[leftmargin=*, itemsep=3mm]

\item Set \texttt{relax\_W = 1} and \texttt{relax\_T = 1}, which takes the full
undamped step. Record what happens: the iteration count, whether it converges,
and any error identifier you get.
\answerspace{4mm}

\item Now set both to \texttt{0.1}. Record the iteration count and the answer.
Compare with the default \texttt{0.5}. Fill this in:

\begin{center}
\begin{tabular}{@{}lccc@{}}
\toprule
$w$ & converged? & passes & $W_{TO}$ [lbf] \\
\midrule
1.0  & & & \\[2mm]
0.5  & & & \\[2mm]
0.1  & & & \\
\bottomrule
\end{tabular}
\end{center}

\item A smaller $w$ took more passes but got there. Explain, in terms of the
slope of the closure map, why damping the step can turn a diverging iteration
into a converging one. Section 8 of the guide gives you the relation you need.
\answerspace{12mm}

\item The \texttt{SizingSteps.relax} docstring says an aircraft whose empty
weight is mostly \emph{fixed} equipment, rather than scaling with size, needs a
smaller $w$. Explain why, using the closure equation.
\answerspace{12mm}

\end{enumerate}

% ------------------------------------------------------------------
\section*{Part B --- Order inside the loop}

\begin{enumerate}[leftmargin=*, itemsep=3mm, resume]

\item Move step 3 (\texttt{prop.T\_SL = T\_SL}) so that it runs \emph{after}
step 4 (the design-point solve) instead of before. Re-run. By how much does
$W_{TO}$ change? Does the checker still pass?
\answerspace{6mm}

\item That change is small, but it is still wrong. Say precisely what the
design-point solve is now reading, and why it is stale. Which physical chain
connects \texttt{prop.T\_SL} to the constraint curves on the F-16?
\answerspace{12mm}

\item Move step 2 (the tail sizing) to \emph{after} step 5 (the weights call).
Re-run. What changes, and which quantity is now one pass out of date?
\answerspace{8mm}

\end{enumerate}

% ------------------------------------------------------------------
\section*{Part C --- When no aircraft closes}

\begin{enumerate}[leftmargin=*, itemsep=3mm, resume]

\item Before the loop, multiply the payload by 3:
\begin{center}
\texttt{wts.W\_payload\_fixed = 3 * wts.W\_payload\_fixed;}
\end{center}
Re-run and record the error identifier and the reported \texttt{denom}.
\answerspace{6mm}

\item State what a non-positive denominator means \emph{physically}. Your
answer should not contain the word ``denominator''.
\answerspace{12mm}

\item Increase the payload gradually and find roughly where the loop stops
closing. Report the largest payload that still converges, and the $W_{TO}$
there. What has happened to the empty and fuel fractions at that point?
\answerspace{10mm}

\end{enumerate}

% ------------------------------------------------------------------
\section*{Part D --- Growth factor}

\begin{enumerate}[leftmargin=*, itemsep=3mm, resume]

\item Size the baseline aircraft, then size it again with 200 lbf more payload.
Compute
\[
\text{growth factor} = \frac{\Delta W_{TO}}{\Delta W_{payload}}.
\]
Report your number. Build a \textbf{fresh stack} for each run.
\answerspace{6mm}

\item Compare it with $1/\mathrm{denom}$ at the baseline. Which is larger?
Section 10 of the guide explains why the two differ for the Boeing 777 --- does
the same explanation apply to the F-16A, or does this aircraft behave the other
way? Look at how \texttt{denom} moves with $W_{TO}$ in your history table.
\answerspace{14mm}

\item The F-16A growth factor is larger than the 777's. Give one reason rooted
in the weight breakdown of a fighter against a transport.
\answerspace{12mm}

\end{enumerate}

% ------------------------------------------------------------------
\section*{Part E --- Level 1 against Level 2}

\begin{enumerate}[leftmargin=*, itemsep=3mm, resume]

\item Run \texttt{LAB00} (Level-1 loop) and your \texttt{LAB01} (Level-2 loop)
and fill in:

\begin{center}
\begin{tabular}{@{}lcc@{}}
\toprule
& Level-1 loop & Level-2 loop \\
\midrule
$W_{TO}$ [lbf]  & & \\[2mm]
$T_{SL}$ [lbf]  & & \\[2mm]
$S_{ref}$ [ft$^2$] & & \\[2mm]
$S_{ht}$ [ft$^2$]  & & \\[2mm]
passes          & & \\
\bottomrule
\end{tabular}
\end{center}

Note that the two use different \emph{discipline} models as well as different
loops, so this is not a clean comparison of the loops alone. Say which
differences you can and cannot attribute to the loop.
\answerspace{12mm}

\item In \texttt{LAB01}, set \texttt{geom.S\_ref = 120} immediately after
building the stack, then run. Do the same in \texttt{LAB00}. Which loop is more
disturbed by the bad starting wing, and why?
\answerspace{12mm}

\end{enumerate}

% ------------------------------------------------------------------
\section*{Part F --- Reading the result}

\begin{enumerate}[leftmargin=*, itemsep=3mm, resume]

\item After a converged run, call
\begin{center}
\texttt{[\textasciitilde, \textasciitilde, info] = con.optimal\_point\_continuous([WS, TW]);}
\end{center}
and look at \texttt{info.active\_names}. Which requirements are binding? What
would happen to the sized aircraft if you relaxed one of them, and what would
happen if you relaxed one that is \emph{not} in the list?
\answerspace{14mm}

\item Your loop converges to roughly 23{,}300 lbf. The Brandt F-16A reference
is 31{,}377 lbf. Give two distinct reasons for the gap that have nothing to do
with whether your loop is correct.
\answerspace{14mm}

\end{enumerate}

\vfill
\begin{warnbox}{A reminder about handles}
The discipline objects are handles, and a sizing run writes into them. If two
of your experiments disagree with each other, the first thing to check is
whether you rebuilt the stack in between.
\end{warnbox}

\end{document}
